\section{Evaluation Details}
\label{app:evaluation-details}
\subsection{Coverage, Completeness Scope, and Unresolved Cases}
\label{app:coverage-and-acceptance}
Of 209 cases, 123 are modeled: 85 require verification only and 38 require
axis inference, of which 31 succeed. The other 86 lack established semantic
support. Acceptance requires proof, compilation, and portable validation
under the written guards and supplied fields. The 31/38 total combines
31/32 successes in the initial run and 0/6 in the slice/convolution run;
a later single-compiler regression reproduces it. One broadcast--transpose
case uses an equivalent supplied-shape calculation; its original frozen
version remains unresolved.

\paperparagraph{Representative rewrites.}
Write $D_a$ for generalized dot with its four batch/contraction lists,
$T_p$ for transpose, and $?a$ for inferred fields. Under the recorded
matrix-shape and LHS-axis conditions, representative cases are
\begin{equation}
\begin{aligned}
 D_a(D_b(X,Y),Z)&\Rightarrow D_{?c}(X,D_{?d}(Y,Z)),
 &&\text{eight dot lists},\\
 T_p(T_qX)&\Rightarrow T_{?r}X,
 &&\text{one permutation}.
\end{aligned}
\end{equation}
The transpose case allows arbitrary rank.

\paperparagraph{Restricted structural extension.}
Theorem~\ref{thm:complete-synthesis} covers the ten case forms in the first
row of Table~\ref{tab:xla-theorem-scope} when every input extent is at least
two and values are independent exact reals. The extension admits positive
extents, including singletons, and supplied-extent broadcast. Operations
connect whole axes by identity or reflection, introduce axes of supplied
size, and sum over full axes. They may multiply operand values and a fixed
nonzero rational weight. Equal sizes alone never identify coordinates.
Each connected class has at most one live axis per operand and one output
position; connections form a forest, with no reflection cycles.
Captures remain read-once within each component, with distinct capture sets
across additive components and independent holes at original RHS nodes.
Scalar factors are independent captures, not inferred weights.
Broadcast over addition uses one shared axis choice for all components.
Checked unit-only reshape preserves the ordered non-unit coordinates for
\emph{every} legal completion; adapter positions stay fixed. General reshape
does not satisfy this condition.

\begin{table}[!htbp]
\caption{Scope of the 38 modeled axis cases. Counts partition cases, not
full XLA domains. The first two rows total 31 structural candidates.}
\label{tab:xla-theorem-scope}
\small
\setlength{\tabcolsep}{4pt}
\begin{tabularx}{\linewidth}{@{}p{0.46\linewidth}rY@{}}\toprule
Case group or case & Count & Scope / recorded result or first failure\\\midrule
Main-theorem forms: dot, transpose, sum & 10
& All inferred; extent $\ge2$, exact reals, D1--D4.\\
Additional structural candidates: broadcast, reverse, products, negation, unit views & 21
& All inferred; positive extents and extension obligations below.\\
\nolinkurl{catalog_broadcast_divide_broadcast} & 1
& Unresolved normal-form equality; nonzero divisor not guarded.\\\midrule
\nolinkurl{slice_reverse} & 1
& Derived metadata not established within its attribute type.\\
\nolinkurl{expr_slice_reverse_general} & 1
& Comprehension outside the pointwise metadata fragment.\\
\nolinkurl{expr_convolution_scale_filter_bias} & 1
& Intermediate shape equality not established.\\
\nolinkurl{catalog_broadcast_slice_broadcast_only} & 1
& Extent inequality in supplied metadata not established.\\
\nolinkurl{catalog_broadcast_slice_broadcast} & 1
& Sliced input bounds not established.\\
\nolinkurl{catalog_dot_slice_output} & 1
& Positive extents for an attribute not established.\\\bottomrule
\end{tabularx}
\end{table}

\paperparagraph{Proof sketch.}
For fixed metadata, a canonical state records the ordered shape, the signed
partition of observable input/output ports, and a rational weight:
\[
 W_v=(\mathbf n_v,\Pi_v,\varepsilon_v,\lambda_v).
\]
Erase unit-coordinate relations while retaining shape. A summed class with
no observable port contributes its extent to $\lambda_v$.
The separating-input argument of Appendix~\ref{app:axis-exactness} then gives
state exactness; positive extents and nonzero weights are essential.
At each RHS node, enumerate all finite legal local axis choices and child
states, preserving every supplied field, and retain one witness per state.
Induction proves that every retained state is realizable and every legal
assignment is represented. State equality preserves parent legality and
meaning, so merging equal states loses no completion. Root-state equality
therefore decides feasibility and recovers a valid assignment. For unary
linear operations around addition, retain joint component states and one
original-node witness to preserve shared choices. This extends the induction
in Appendix~\ref{app:correctness-and-completeness}.

\paperparagraph{What the count establishes.}
The 31 candidates pass the structural checks, including capture use,
distinct holes, and three unit-view cases. The theorem and extension apply
only under their stated assumptions, with defined supplied programs and
normalization valid for every completion. The recorded guards do not ensure
these conditions over the full XLA domains, which may include zero extents
and machine arithmetic. The count therefore establishes conditional
structural coverage. The production solver does not fully implement the
exhaustive state-table algorithm.

The table lists the seven unresolved cases and their first failures.
The six slice/convolution cases need further proved adapters to enter either
contract. Scalar division requires an independent nonzero divisor under
exact-real arithmetic, but the measured rule lacks this guard. These failures
do not prove impossibility, and resolving them may expose further obligations.

\FloatBarrier
\subsection{Automatic Attribute-Inference Scalability}
\label{app:automatic-attribute-inference-scalability}
\paperfigure{1}{figures/scalability.pdf}
{Automatic attribute-inference scalability. Compilation time by phase for
(a) dot-chain length and (b) additive-term count, and (c) runtime metadata
execution time by phase as tensor rank increases. Markers show mean total
time; whiskers show the minimum and maximum.}
{fig:scalability}
{Three panels show compilation costs for dot chains and additive terms up
to 64, and runtime feasibility-checking and attribute-construction costs
for tensor ranks up to 128. Stacked bars show mean phase costs.}

We evaluate automatic attribute-inference scalability by separating
compilation of the feasibility check $q$ and attribute program $f$ from
their runtime execution (Figure~\ref{fig:scalability}). All runs use an
Intel Core Ultra 7 155H CPU, with CPU-only, single-threaded execution.
Panels (a) and (b) use chains of arbitrary-rank dot operations with all RHS
attributes inferred automatically. Panel (a) reassociates a left-associated
chain into a right-associated chain, increasing dots per side through
1, 2, 4, 8, 16, 32, and 64. Panel (b) varies the number of independent
additive terms over the same sizes, with two dots per term and balanced
addition trees. All tested sizes compile successfully. Mean compilation
time grows approximately linearly, reaching about 39\,ms for 64 dots per
side and 60\,ms for 64 additive terms. Panel (c) compiles one rank-independent
two-dot rule once, then varies the concrete tensor rank. Feasibility
checking and attribute construction are directly timed within the same
execution. Total runtime metadata execution grows with rank, reaching
about 451\,$\mu$s at rank 128; rank 256 exceeds the metadata work budget.
Thus, TEPL compiles the tested rank-polymorphic rules in tens of
milliseconds at the largest sizes, and a single compilation handles
varying tensor ranks without rank-specific recompilation.
